\chapter{A virtual laboratory for the dimension witness}
\label{ch:dimension-lab}

The finite-record dimension test has a useful separation between validity and
power. Fixed preparation error, detector confusion, and independent noise on
each repeated step can substantially weaken a quantum signal while preserving
the classical null theorem. A small coherent frequency error is particularly
costly: it moves the recurrence away from the prescribed observation grid.
An independently calibrated \emph{common} clock correction repairs this failure
without changing the classical dimension bound. The following experiments
establish that statement within an independently implemented simulator. They
are not observations from a physical device.

\section{Physical simulation and frozen experiments}
The base system is a driven qubit with angular drive frequency $\omega=1$ and
isotropic depolarization rate $\gamma=10^{-3}$, in declared model time units.
The observation step is $\Delta=2\pi/7$. A fixed erroneous preparation places
probability $e_{\rm prep}$ on $|1\rangle$. A fixed binary detector has false-positive
and false-negative probabilities $e_0,e_1$. The drive can have a fixed amplitude
error, fixed detuning, and independent two-point amplitude and duration noise
refreshed on every step.

The primary simulator exponentiates the three-dimensional Bloch generator,
averages its four independently refreshed amplitude/duration channels, and
then propagates the state. It does not call the dimension certificate's
sinusoidal probability helper. A separate four-dimensional vectorized Lindblad
calculation reproduces the probabilities to within $1.45\times10^{-15}$ across
the initial 24 declared regimes. One additional record stream simulates the
preparation bit, every step's noise choice, final Born measurement, and detector
confusion explicitly. Its $18\,300$ records and noise history are retained in
\texttt{dimension\_lab\_shots.npz}.

The initial contract was saved before generating its pilot or production data.
It fixes all settings, seeds, noise regimes, and the original $18\,300$-record
critical-count test. Each power estimate uses $20\,000$ independently generated
pseudoexperiments. A pilot of $20\,000$ records at each of 13 settings selects
an alternative critical-count plan from a frozen candidate family. Its
$260\,000$ additional records are separate from the production data. Later
clock-repair and total-budget comparisons have separate frozen contracts and
new seeds. All declared cases, including failed regimes, are retained.

The saved shot ordering is a random permutation of independent simulated
records. It represents randomized acquisition under this simulator; it is not
evidence that an actual acquisition was stable or robust against drift.

\section{Which apparatus errors preserve the null}
Let $\alpha$, $P$, and $f$ denote the actual prepared distribution, repeated
classical channel, and stochastic binary detector. The null theorem concerns
$p_k=\alpha P^k f$, regardless of whether these quantities match nominal values.
Consequently a fixed imperfect preparation or readout is already covered.
Independent identical step errors replace $P$ by their average channel and
are also covered. No assumed qubit frequency, pure preparation, or detector
fidelity enters this classical argument.

An error drawn once and held for an entire run is different. Its record is
generally $\mathbb E[\alpha P(V)^k f]$, not
$\alpha(\mathbb E P(V))^k f$. The retained noise label carries memory and must
be included in a classical state count. In the declared quasi-static control,
these two probability sequences differ by as much as $0.20164$. Likewise,
allowing the detector response itself to depend on $k$ lets one frozen
classical state reproduce the full target sequence. The resulting rejection
rate is high, as intended for this outside-contract control. A dimension
claim therefore needs a common interface and channel, with any controller
memory included.

\section{Noise tolerance and a matched production comparison}
The original test's simulated power is $98.565\%$ in the nominal model.
Representative single-error regimes show the following behavior:
\begin{center}
\begin{tabular}{lr}
Illustrative error & Original-test power\\
\hline
Detector false-positive probability $0.0005$ & $95.245\%$\\
Detector false-positive probability $0.001$ & $81.260\%$\\
Detector false-positive probability $0.002$ & $28.050\%$\\
Preparation error probability $0.001$ & $81.330\%$\\
Independent step amplitude jitter $\pm1\%$ & $98.145\%$\\
Independent step duration jitter $\pm1\%$ & $98.190\%$\\
Fixed amplitude offset $+0.5\%$ & $63.920\%$\\
Fixed amplitude offset $+1\%$ & $1.595\%$
\end{tabular}
\end{center}
These are illustrative dimensionless settings, not a claim about the calibration
or feasibility of any particular hardware platform. Complete marginal
Clopper--Pearson intervals appear in the machine-readable results. Monte Carlo
uncertainty and unknown real-device systematics are separate matters.

The combined design regime has preparation error $0.0005$, detector errors
$(e_0,e_1)=(0.0005,0.003)$, amplitude offset $0.003$, detuning $0.02$, independent
amplitude jitter $\pm0.005$, and duration jitter $\pm0.002$. Its original-test
power is $67.66\%$, with a $95\%$ Monte Carlo interval
$[67.01\%,68.31\%]$. The independent pilot selects
\[
 \epsilon_0=61/10000,\qquad \beta_0=1503/10000,
 \qquad \tau=781/10000.
\]
The resulting critical counts are
\[
 (314,22,22,22,22,298,49,298,22,22,22,22,282),
\]
with the low-count inequality at setting seven and high-count inequalities
elsewhere. Exact rational tail checks give a type-I bound below $0.049144$.
At the unchanged production allocation it obtains $71.28\%$ power, with
interval $[70.65\%,71.91\%]$. The paired improvement is $3.62$ percentage
points; a conservative Hoeffding interval is $[1.70,5.54]$ points.
This is a production-budget comparison conditional on an additional pilot,
not a total-cost superiority claim.

A library of 400 homogeneous classical models, spanning one through four
states and both continuous-time and discrete stochastic channels, supplies
12 adversarially selected null examples. None rejects in $20\,000$ repetitions
per selected example. A zero count has marginal upper $95\%$ limit
$0.018443\%$. These finite checks are supplementary: the exact type-I
certificate, rather than this library, establishes validity over the full null
class.

\section{Repairing coherent calibration error with one common clock}
The fixed-amplitude failure is not merely a larger diagonal leakage. At a
$1\%$ amplitude error the last setting has mean probability approximately
$0.148983$, below the original required off-diagonal boundary $0.15$.
Averaging this particular criterion more accurately will not repair that
mismatch. The observation grid itself needs attention.

In a second independently frozen loop, a $260\,000$-record pilot fits one drive
amplitude offset with the nominal Bloch model. The production step is then
multiplied by $1/(1+\widehat a)$. All 13 observation times use this \emph{same}
step; none is individually corrected. The fit proposes a design and can be
wrong. It does not supply a premise for the null theorem.

\begin{proposition}[Independent common-clock calibration preserves validity]
Condition on arbitrary pilot data $D$, and choose a positive common step
$\Delta(D)$ and a production plan whose fixed-design null error is at most
$\alpha$. If the independent production records share one actual preparation,
channel, and detector, and every production setting repeats that chosen step,
then their unconditional null error is at most $\alpha$.
\end{proposition}
\begin{proof}
For each fixed pilot outcome, the production means factor through one actual
channel $P_D$ as $p_k=\alpha_D P_D^k f_D$. Apply the fixed-design dimension
certificate and count-tail guarantee conditionally on $D$, and average the
resulting bound over $D$. No accuracy condition on the frequency fit is used.
\end{proof}

The second holdout uses the original $18\,300$ record quotas and original
critical counts. Its outcomes are:
\begin{center}
\begin{tabular}{lrrr}
Regime & Original clock & Common-clock repair & Repair's $95\%$ interval\\
\hline
Amplitude $+0.5\%$ & $63.740\%$ & $98.610\%$ & $[98.44,98.77]\%$\\
Amplitude $+1\%$ & $1.575\%$ & $98.720\%$ & $[98.55,98.87]\%$\\
Amplitude $+2\%$ & $0\%$ & $98.595\%$ & $[98.42,98.75]\%$\\
Combined design noise & $67.855\%$ & $80.380\%$ & $[79.82,80.93]\%$
\end{tabular}
\end{center}
The four positive fitted offsets shorten the production evolution time by
approximately $0.30\%$ to $1.99\%$. The shot count and $126\,900$ repeated
channel applications remain identical. Reset, readout, and controller overhead
are not included in this evolution-time accounting.

\section{Calibration is a resource, and sometimes more averaging wins}
Each frequency pilot costs $260\,000$ records and $1\,820\,000$ channel
applications. Combined with one production experiment, its total is
$278\,300$ records. If one calibration remains valid for 100 later production
experiments, its record overhead amortizes to about $14.21\%$ per experiment;
for 1000 uses, about $1.42\%$. Such reuse assumes drive stability and is not
established by the simulator.

An additional frozen comparison spends all $278\,300$ records on the original
clock, proportionally enlarges the original quotas, and recalibrates the
original forbidden-region test at the same $5\%$ type-I level. Its measured
powers are $100\%$, $1.31\%$, $0\%$, and $100\%$ for the four respective
regimes in the preceding table. Thus extra averaging wins the mild-offset and
combined-noise cases at this total record budget. Common retiming repairs the
two severe phase failures. The equal-record baseline retains the original
probability region and is not an oracle or an optimal competing protocol.
The total evolution durations differ and are included in the result files.
These comparisons support choosing an intervention according to the actual
failure mechanism; they do not establish universal cost superiority.

\section{Reproduction and interpretation}
The module \texttt{dimension\_lab.py} exposes the physical channel simulators,
exact plan checks, pilot design, record generation, and three frozen runs.
The contracts, selected thresholds, all noise-regime results, count intervals,
raw shot stream, and seeds are retained. Twelve producer tests include an
independent density-matrix solver, direct measurement-noise checks, timing-memory
counterexamples, exact null checks, malformed-plan rejection, reproducible shot
streams, and common-clock repair.

The concrete result is an experimentally relevant failure and repair:
coherent grid error can dominate the power of a valid dimension witness,
and independent calibration of one shared clock can restore power while
preserving its target-independent null guarantee. Physical device validation,
external review, and any claim of historical originality remain separate work.
